% ============================================================
%  USPSC-Cap4-Avaliacao_Experimental_4.4-Discussao.tex
%  Seção 4.4 — Discussão
% ============================================================

\section{Discussão}
\label{sec:discussao_resultados}

\subsection{Especialização Prévia versus Ajuste Local ao Gênero}
\label{subsec:discussao_especializacao}

O contraste mais relevante que emerge dos resultados é o que opõe o FinBERT-PT-BR ao BERTimbau ajustado, duas abordagens neurais que diferem fundamentalmente na estratégia de especialização ao domínio e ao gênero. O FinBERT-PT-BR incorpora \textbf{especialização prévia de domínio} (\textit{prior domain specialization}): o modelo foi pré-treinado sobre um corpus de notícias financeiras formais em português \cite{santos2023finbert}, o que lhe confere vocabulário financeiro robusto, mas não o prepara para as particularidades do gênero textual de publicações do X/Twitter --- sentenças curtas, uso de emojis convertidos, menções a tickers, estilo telegráfico e ausência de marcadores explícitos de polaridade. O BERTimbau ajustado incorpora \textbf{especialização pós-hoc de gênero e subdomínio} (\textit{post hoc genre and subdomain specialization}): o modelo base generalista \cite{souza2020} foi adaptado diretamente sobre publicações do mesmo gênero e subdomínio que compõem o corpus de avaliação.

Os resultados confirmam a hipótese de que a adaptação local ao gênero produz ganhos substanciais mesmo quando comparada a um modelo com especialização de domínio prévia: o BERTimbau ajustado superou o FinBERT-PT-BR em 34,2 pontos percentuais de F1-score macro (77,45\% versus 43,20\%). Esse achado está alinhado com a tradição de \textit{fine-tuning} de modelos de linguagem em domínios especializados \cite{devlin2019,santos2023finbert}, que demonstra que a proximidade entre os dados de ajuste e os dados de inferência é um fator determinante para o desempenho em tarefas de classificação de sentimento.

\subsection{Diagnóstico do Viés do FinBERT-PT-BR}
\label{subsec:discussao_vies_finbert}

O viés do FinBERT-PT-BR em favor da classe negativo (revocação de 88\%, precisão de 20\%) merece análise específica. Uma hipótese explicativa plausível é a distribuição assimétrica de sentimento no corpus de treinamento do modelo: notícias financeiras formais tendem a tratar de eventos de risco, quedas de indicadores, inadimplência e crises, gerando uma representação mais densa de vocabulário negativo no espaço de \textit{embeddings} do modelo. Quando aplicado a publicações do X/Twitter, o modelo interpreta o vocabulário neutro de cobertura de mercado (e.g., ``variação'', ``fechamento'', ``resultado'') como sinal negativo, produzindo o padrão de superclassificação observado. Esse diagnóstico é consistente com a observação de \citeonline{loughran2011}, que documentaram que léxicos de sentimento geral aplicados ao domínio financeiro frequentemente atribuem polaridade incorreta a termos com significado neutro ou positivo em contextos de negócios.

\subsection{Limitações das Abordagens Léxicas em Corpus de Tweet}
\label{subsec:discussao_lexical}

O desempenho das abordagens léxicas --- OpLexicon (F1 macro 26,83\%) e SentiLex-PT (F1 macro 30,79\%) --- é sistematicamente inferior ao que seria esperado de uma aplicação ao domínio para o qual foram construídos: textos em português com marcadores explícitos de polaridade. Dois fatores estruturais explicam esse resultado. Primeiro, publicações jornalísticas financeiras do X/Twitter são predominantemente informativas e factuais; a proporção de termos com polaridade lexical explícita é baixa, o que reduz a capacidade de discriminação dos léxicos. Segundo, a ausência de mecanismos de desambiguação contextual implica que termos polissêmicos ou dependentes de contexto (e.g., ``alta'' como variação de preço versus expressão de intensidade) são classificados com base apenas em sua ocorrência, independentemente do sentido contextual \cite{liu2020survey}.

Ressalta-se que o OpLexicon e o SentiLex-PT foram construídos para o português de domínio geral, sem especialização para o vocabulário financeiro; o desempenho observado pode ser interpretado como um \textit{baseline} léxico de referência, não como o limite superior das abordagens baseadas em léxico para o domínio.


\subsection{Restrições à Generalização dos Resultados}
\label{subsec:discussao_restricoes}

Os resultados devem ser interpretados considerando as restrições metodológicas identificadas na \autoref{sec:limitacoes}. O conjunto de teste foi anotado por um único avaliador, inviabilizando o cálculo de concordância interanotadores; a margem de incerteza associada à qualidade do \textit{gold standard} não é quantificada. Adicionalmente, o corpus é restrito a dois veículos de mídia financeira institucional (InfoMoney e InvestingBrasil), o que circunscreve a representatividade dos resultados ao gênero editorial financeiro em língua portuguesa --- os achados podem não ser diretamente generalizáveis a outros gêneros textuais do domínio financeiro, como análises de analistas independentes, fóruns de investidores ou mensagens de corretoras.

Por fim, o tamanho do corpus de ajuste do BERTimbau (aproximadamente 2.600 instâncias de treinamento) é modesto em comparação com os corpora utilizados em estudos de \textit{fine-tuning} de grande escala; a utilização de conjuntos de dados maiores e de maior diversidade de gênero poderia produzir modelos com desempenho superior nas classes minoritárias, especialmente positivo (F1 = 0,72) e negativo (F1 = 0,69).